\chapter{The Big Picture: Eight Stages and a Loop}
\label{ch:bigpicture}

\Cref{ch:problem} argued the \emph{why}: hallucination is a structural property
of stateless, overconfident, static generation, and the only honest response is
a system that attacks all three at once. This chapter gives the \emph{shape} of
that system in a single sitting. We will not prove anything here and we will not
open any one component. The goal is a reliable mental map, so that when a later
chapter spends thirty pages on the verifier or the router, we already know where
that piece sits, what feeds it, and what it feeds.

\begin{intuition}[The whole machine in three sentences]
\caem{} runs every question through a short assembly line that decides how hard
to work, produces an answer, and then judges that answer before deciding whether
to trust it, refuse, or file it away for later. That assembly line is the
\textbf{per-query pipeline}, and it runs thousands of times without ever changing
the model. Wrapped around it is a slower \textbf{cycle loop} that periodically
retrains the model on the answers the pipeline judged trustworthy, with a guard
that undoes any retraining that makes the model worse. One fast clock answers
questions; one slow clock makes the answerer better. Everything in the book is a
detail of one of those two clocks.
\end{intuition}

\section{Two clocks, not one}
\label{sec:two-clocks}

The single most important idea in this chapter is that \caem{} has \emph{two
independent clocks}, and almost every later confusion dissolves once we keep them
apart.

\begin{description}[leftmargin=0pt, labelsep=0.5em, style=nextline]
  \item[The per-query clock (Stages 1 to 7).] One tick is one question. The
  system encodes the query, decides which of three tiers should handle it,
  generates or retrieves an answer, scores that answer with a verifier, and
  applies a four-way decision. This clock \emph{reads} the current memory, the
  current model weights, and the current calibration, and at most it \emph{writes
  one episode} into memory. It never trains anything. It ticks thousands of times
  per cycle.

  \item[The per-cycle clock (Stage 8).] One tick is one \emph{cycle boundary},
  reached after a batch of queries has been processed. At a boundary the system
  fine-tunes the model on its own verified answers, re-checks that it has not
  forgotten anything, recalibrates the verifier, and re-scores the stored
  memory. This clock \emph{mutates} the weights, the calibration, and the
  consolidated memory. It ticks a handful of times across the whole run, six
  times in the realised trajectory of \Cref{ch:trajectory}.
\end{description}

The two clocks are nested: the per-cycle clock is the outer loop, and each of its
ticks drives the per-query clock thousands of times in batch. Stages 1 through 7
are where the model, the calibration, and the memory are merely \emph{consumed};
Stage 8 is the only place they are \emph{changed}. \Cref{fig:two-clocks} draws
the nesting; the rest of the chapter walks the inner clock first, then the outer.

\begin{figure}[htbp]
\centering
\begin{tikzpicture}[
  font=\small\sffamily,
  outer/.style={rectangle, rounded corners=8pt, draw=cbFormula!70!black,
    fill=cbFormula!5, line width=1.2pt},
  inner/.style={rectangle, rounded corners=6pt, draw=cbIntuition!75!black,
    fill=cbIntuition!8, line width=1pt},
  cbstep/.style={rectangle, rounded corners=3pt, draw=cbFormula!70!black,
    fill=cbFormula!12, line width=0.8pt, align=center, inner sep=4pt,
    font=\footnotesize, minimum height=0.8cm},
  flow/.style={-{Stealth[length=2.2mm]}, line width=1pt, draw=cbInk!70},
  carry/.style={-{Stealth[length=2.2mm]}, line width=1pt, draw=cbExample!75,
    dashed},
]
  % outer frame
  \node[outer, minimum width=14.5cm, minimum height=5.6cm] (O) at (0,0) {};
  \node[anchor=north west, font=\footnotesize\bfseries\sffamily,
        text=cbFormula!75!black] at (O.north west)
        {\,\,PER-CYCLE CLOCK \;(Stage 8): ticks a handful of times};

  % inner frame (the per-query pipeline)
  \node[inner, minimum width=9.2cm, minimum height=2.5cm] (I) at (-2.5,0.45) {};
  \node[anchor=north west, font=\footnotesize\bfseries\sffamily,
        text=cbIntuition!80!black] at (I.north west)
        {\,\,PER-QUERY CLOCK \;(Stages 1--7): ticks thousands of times};
  \node[align=center, font=\footnotesize, text=cbInk!85] at (-2.5,0.2)
    {one question in, one answer out\\[1pt]
     \emph{reads} weights, memory, calibration\\[1pt]
     \emph{writes} at most one episode};

  % cycle-boundary steps on the right
  \node[cbstep, text width=3.8cm] (sil)  at (5.0, 1.8) {Fine-tune on verified\\answers + retention guard};
  \node[cbstep, text width=3.8cm] (recal)at (5.0, 0.45) {Recalibrate verifier\\+ retro-verify memory};
  \node[cbstep, text width=3.8cm] (grow) at (5.0,-0.9) {Grow memory with\\fresh-data stream};

  % flows
  \draw[flow] (I.east) -- ++(0.5,0) |- (sil.west);
  \draw[flow] (sil) -- (recal);
  \draw[flow] (recal) -- (grow);
  \draw[carry] (grow.south) |- ($(I.south)+(0,-0.95)$)
        node[pos=0.55, below, font=\scriptsize, text=cbExample!60!black, text width=7cm, align=center]
        {updated weights, memory, calibration carried into next cycle}
        -- (I.south);
\end{tikzpicture}
\caption[The two clocks]{The two clocks of \caem{}. The inner per-query pipeline
(Stages 1 to 7) answers questions without ever training. The outer per-cycle loop
(Stage 8) is the only place the weights, calibration, and consolidated memory
change. Each cycle boundary fine-tunes the model on the answers the pipeline
judged trustworthy, then hands the updated state back to the inner clock. (Schematic.)}
\label{fig:two-clocks}
\end{figure}

\section{The per-query pipeline at a glance}
\label{sec:pipeline-glance}

A single question enters through one public entry point and leaves as one
structured record. Everything between is a fixed sequence of stages.
\Cref{fig:master-pipeline} is the master diagram of the book; it is worth a long
look, because every chapter in Parts II and III is a zoom into one of its boxes.

\begin{figure}[htbp]
\centering
\resizebox{\textwidth}{!}{%
\begin{tikzpicture}[
  font=\small\sffamily,
  io/.style={rectangle, rounded corners=2pt, draw=cbInk!55, fill=cbInk!6,
    line width=1pt, align=center, inner sep=5pt, minimum height=1.0cm},
  front/.style={rectangle, rounded corners=4pt, draw=cbIntuition!75!black,
    fill=cbIntuition!10, line width=1pt, align=center, inner sep=5pt,
    minimum height=1.1cm, drop shadow={shadow xshift=0.3ex,
    shadow yshift=-0.3ex, fill=black!10}},
  mem/.style={rectangle, rounded corners=4pt, draw=cbExample!70!black,
    fill=cbExample!12, line width=1pt, align=center, inner sep=5pt,
    minimum height=1.1cm, drop shadow={shadow xshift=0.3ex,
    shadow yshift=-0.3ex, fill=black!10}},
  tier/.style={rectangle, rounded corners=4pt, draw=cbScenario!70!black,
    fill=cbScenario!10, line width=1pt, align=center, inner sep=4pt,
    minimum height=1.0cm, text width=2.4cm, drop shadow={shadow xshift=0.3ex,
    shadow yshift=-0.3ex, fill=black!10}},
  judge/.style={rectangle, rounded corners=4pt, draw=cbFormula!70!black,
    fill=cbFormula!10, line width=1pt, align=center, inner sep=5pt,
    minimum height=1.1cm, drop shadow={shadow xshift=0.3ex,
    shadow yshift=-0.3ex, fill=black!10}},
  flow/.style={-{Stealth[length=2.6mm]}, line width=1.1pt, draw=cbInk!75},
  bypass/.style={-{Stealth[length=2.6mm]}, line width=1.1pt, draw=cbScenario!80,
    dashed},
  store/.style={-{Stealth[length=2.6mm]}, line width=1.1pt, draw=cbExample!75,
    dashed},
]
  % top: memory store, spanning
  \node[mem, minimum width=4.2cm] (store) at (6.0, 3.0)
    {\textbf{Episodic memory}\\{\footnotesize verified episodes (FAISS)}};

  % front end
  \node[io] (q) at (-1.0,0) {Query};
  \node[front, text width=2.6cm] (enc) at (2.0,0)
    {\textbf{Encode + $\upre$}\\{\footnotesize Stages 2, 1}};
  \node[front, text width=1.9cm] (route) at (5.3,0)
    {\textbf{Router}\\{\footnotesize Stage 3 / 3b}};

  % three tiers
  \node[tier] (t1) at (8.6, 1.5) {\textbf{Tier 1}\\memory-direct};
  \node[tier] (t2) at (8.6, 0.0) {\textbf{Tier 2}\\zero-shot};
  \node[tier] (t3) at (8.6,-1.5) {\textbf{Tier 3}\\RAG};

  % verify + gate + output
  \node[judge, text width=1.9cm] (ver) at (11.7,-0.75)
    {\textbf{Verify}\\{\footnotesize Stage 5}};
  \node[judge, text width=2.0cm] (gate) at (14.6,-0.75)
    {\textbf{Gate}\\{\footnotesize Stage 6}};
  \node[io, text width=2.0cm] (out) at (14.6, 1.5)
    {\textbf{Output}\\{\footnotesize Stage 7}};

  % front-end flow
  \draw[flow] (q) -- (enc);
  \draw[flow] (enc) -- (route);
  \draw[flow] (enc.north) -- ++(0,1.0) -| (store.west)
    node[pos=0.25, above, font=\scriptsize, text=cbInk!60] {probe};

  % route to tiers
  \draw[flow] (route.east) -- ++(0.5,0) |- (t1.west);
  \draw[flow] (route.east) -- ++(0.5,0) |- (t2.west);
  \draw[flow] (route.east) -- ++(0.5,0) |- (t3.west);

  % tiers to verify (T2/T3) and bypass (T1)
  \draw[flow] (t2.east) -- ++(0.4,0) |- (ver.west);
  \draw[flow] (t3.east) -- ++(0.4,0) |- (ver.west);
  \draw[bypass] (t1.east) -- node[above, font=\scriptsize, text=cbScenario!70!black]
    {no verify, no store} (out.west);

  % verify -> gate -> output
  \draw[flow] (ver) -- (gate);
  \draw[flow] (gate.north) -- (out.south);

  % gate store back to memory
  \draw[store] (gate.south) -- ++(0,-1.0) -| (store.west)
    node[pos=0.2, below, font=\scriptsize, text=cbExample!55!black]
    {store / defer};
\end{tikzpicture}}
\caption[The master per-query pipeline]{The per-query pipeline. A query is
encoded and scored for confidence, the router picks one of three tiers, and the
chosen tier produces an answer. Tier 1 (a memory hit) returns its stored answer
directly, skipping verification and storage entirely. Tiers 2 and 3 send their
generated answer through the nine-signal verifier and the four-outcome gate,
which may return it, refuse it, or write it back into memory. The dashed green
path is the only way an episode enters memory; the dashed teal path is the only
tier that bypasses the judge. (Schematic.)}
\label{fig:master-pipeline}
\end{figure}

\subsection*{The seven inner stages in one line each}

Reading \Cref{fig:master-pipeline} left to right, with the chapter that opens each
box:

\begin{enumerate}[leftmargin=1.8em]
  \item \textbf{Pre-routing confidence} ($\upre$, \Cref{ch:upre}). A single short
  forward pass, with no full generation, produces a calibrated guess at how
  likely the model is to answer this query correctly on its own. It fuses a
  token-probability signal with an encoder-convergence signal.
  \item \textbf{Encode and probe} (\Cref{ch:memory}). The query is turned into a
  768-dimensional embedding, which is used both to search the episodic memory for
  the single nearest stored episode and, later, to ground the verifier.
  \item \textbf{Route} (\Cref{ch:router}). The retrieved similarity and the
  stored answer's quality are fused into one dispatch score that selects a tier,
  with a safety override that escalates whenever $\upre$ falls below a fixed
  floor.
  \item \textbf{Refine with the retrieval-utility classifier} (Stage 3b,
  \Cref{ch:ruc}). When the router would default to retrieval, a small learned
  classifier decides whether retrieval will actually help this query, or whether
  the model is better off answering zero-shot. This is the centrepiece of the
  second version of the system.
  \item \textbf{Execute the tier} (\Cref{ch:tiers}). Tier 1 serves the stored
  answer; Tier 2 generates zero-shot on the fine-tuned model; Tier 3 retrieves
  passages from a large public corpus and generates a grounded answer.
  \item \textbf{Verify} (\Cref{ch:verifier}). Every generated answer, but never a
  Tier 1 hit, is scored on nine signals across four families: token and dropout
  uncertainty, internal-state probes, semantic consistency across resampled
  chains, and entailment against retrieved evidence.
  \item \textbf{Decide and output} (\Cref{ch:composite,ch:output}). The nine
  signals are folded into one calibrated probability, the four-outcome gate maps
  that probability to \emph{store}, \emph{defer}, \emph{abstain}, or
  \emph{discard}, and the output stage returns the answer or a calibrated
  refusal.
\end{enumerate}

\begin{fromcode}[Where this lives, and one honest subtlety]
The entire inner pipeline is one method, \code{CAEMPipeline.answer(query,
store\_to\_memory, source\_benchmark)} in \code{caem/pipeline.py}. It is the
single public per-query entry point and it returns one \code{PipelineResult}
object. Two details are worth flagging now so they do not surprise us later.
First, the stage \emph{numbers} are a conceptual taxonomy, not the execution
order: in the actual code the query is encoded, then scored for $\upre$, then the
memory is searched, then the route is chosen. We draw the diagram in code order.
Second, an older design carried a separate post-generation confidence gate; it
was retired, and the verifier of \Cref{ch:verifier} is now the single source of
truth after generation. Where the code comments still mention that retired gate,
they are stale.
\end{fromcode}

\section{The tier asymmetry, the load-bearing branch}
\label{sec:tier-asymmetry}

If only one fact from this chapter survives, it should be this one. The three
tiers are not symmetric, and the asymmetry is the whole point of the design.

\begin{itemize}[leftmargin=1.4em]
  \item \textbf{Tier 1} is a memory hit. The router found a stored episode close
  enough to the query that its verified answer can be served directly. There is
  \emph{no generation} and \emph{no verification}: the answer and its
  already-known trust score are read straight off the stored entry, and the
  pipeline returns in well under half a second. A Tier 1 answer was already
  judged trustworthy on the cycle it was stored, so re-judging it would be wasted
  work.
  \item \textbf{Tier 2 and Tier 3} are fresh generations, zero-shot and
  retrieval-augmented respectively. A fresh generation has never been judged, so
  it \emph{must} pass through the verifier and the gate before the system will
  trust it or store it.
\end{itemize}

This is the branch that turns the three tiers into a self-organising compute
ladder. Cheap, repeated questions collapse onto Tier 1 over time as their answers
get memorised; novel or hard questions climb to Tiers 2 and 3 and pay the full
verification cost. The cost of trust is paid once per fact, not once per ask.

\begin{pitfall}[Tier 1 is not ``the model is sure''; it is ``we have seen this answered well before'']
It is tempting to read the tiers as a confidence ladder, with Tier 1 meaning
``the model is most confident.'' That is wrong. Tier 1 means a \emph{stored
episode} matches the query closely enough, and that episode carries the trust
score it earned when it was first verified. The confidence that the \emph{current}
query is well served comes from the match quality and the stored score together,
not from a fresh judgement of the answer. This is exactly why Tier 1 can skip the
verifier: the judging already happened, on a past cycle, and its verdict was
written into the episode.
\end{pitfall}

One escalation edge is worth noting on the diagram. If a Tier 2 generation comes
back empty or fails, the pipeline quietly escalates that single query to Tier 3
and marks the result as escalated. The reported tier still reads two, but the
answer came from retrieval. It is a small robustness valve, not a separate stage.

\section{A query, end to end}
\label{sec:end-to-end}

Abstractions settle once we watch a real query move through the machine. Two
scenarios follow, the first a claim that ends up stored, the second a question
that ends up refused. The refusal scenario is an actual record from a cold-start
pipeline pass, reproduced with its logged signals. The stored-claim scenario is an
illustrative walkthrough of the defer-then-promote path, hand-authored to make that
two-cycle behaviour concrete rather than transcribed verbatim from a single record.
Together they exercise the two behaviours that define the architecture: promoting
a tentatively answered claim to memory after later review, and turning a
low-confidence question into an honest refusal rather than a fabricated answer.

\begin{scenario}[A claim that gets stored: ``The French Revolution began before 1792'' (FEVER)]
The task is to label the claim as \emph{supports}, \emph{refutes}, or \emph{not
enough information}.
\begin{enumerate}[leftmargin=1.5em, itemsep=2pt]
  \item \textbf{Pre-routing.} The fused, temperature-scaled confidence
  comfortably clears the safety floor, so no override fires.
  \item \textbf{Encode and probe.} On this cycle-zero pass the memory is empty,
  so there is no neighbour to retrieve.
  \item \textbf{Route.} With nothing in memory, the router defaults the query
  toward the retrieval tier.
  \item \textbf{Refine (RUC).} The retrieval-utility classifier overrides that
  default to the \emph{zero-shot} tier, judging the claim short, self-contained,
  and unlikely to benefit from passage retrieval.
  \item \textbf{Execute.} Tier 2 generates the label \emph{supports}.
  \item \textbf{Verify.} The nine signals read as follows: top-passage entailment
  is high (above $0.94$), the pooled-mean and atomic entailment are lower because
  only part of the supporting evidence sits in the context window, semantic
  consistency across resampled chains is high (above $0.85$), and
  question-answer relevance is moderate (around $0.69$). Folded together, the
  calibrated composite is about $0.51$.
  \item \textbf{Decide.} A composite of $0.51$ clears the deferral threshold
  ($\taudefer = 0.45$) but falls just below the storage threshold
  ($\taustore = 0.60$), so the entry is parked in the \emph{deferred buffer} for
  reconsideration at the next cycle boundary, rather than stored or discarded.
  \item \textbf{Cycle boundary (Stage 8).} After the model is fine-tuned, the
  retroactive re-verification pass re-scores the same entry under the improved
  model. This time it crosses $\taustore$ and is promoted into the cold-start
  memory as entry identifier zero.
  \item \textbf{Output.} The user receives the label \emph{supports}, which is
  correct.
\end{enumerate}
\textbf{The lesson:} a borderline-but-promising answer is neither trusted blindly
nor thrown away. It waits, and a later, better model decides its fate.
\end{scenario}

\begin{scenario}[A question that gets refused: ``\,`If I Ruled The World' comes from which musical?'' (TriviaQA)]
\begin{enumerate}[leftmargin=1.5em, itemsep=2pt]
  \item \textbf{Pre-routing.} The pre-routing confidence lands around $0.66$,
  comfortably above the safety floor, so no override fires. On this cold-start
  pass the episodic memory is still empty, so no cheap memory path can be earned
  and the query is scored normally.
  \item \textbf{Refine (RUC).} Consulted on the same query, the classifier
  confirms the retrieval tier: the base model's self-knowledge probe reads low,
  and the dense passage corpus does carry supporting evidence.
  \item \textbf{Execute.} Tier 3 retrieves passages and runs a grounded
  generation, which returns an \emph{empty} answer, an explicit refusal signal.
  \item \textbf{Verify.} The signals on the empty answer are damning:
  token-level uncertainty is very low (around $0.02$), top-passage entailment is
  low (around $0.11$), atomic entailment is low (around $0.09$), and
  question-answer relevance sits at chance (around $0.50$). The calibrated
  composite is about $0.13$, well below both thresholds.
  \item \textbf{Decide.} The composite fails both the storage and deferral
  cut-points, and the top-passage entailment falls below the registered
  abstention floor, so the gate returns \emph{abstain} rather than a silent
  discard.
  \item \textbf{Output.} The user sees the registered refusal string, and the raw
  model prediction is suppressed from the user surface.
\end{enumerate}
\textbf{The lesson:} when the evidence does not support an answer, the system says
so, instead of emitting a confident guess. This is the calibration mismatch of
\Cref{ch:problem} being repaired in practice.
\end{scenario}

\section{What the pipeline hands back}
\label{sec:pipeline-result}

Each tick of the per-query clock returns one structured record, and it is worth
knowing its shape, because the rest of the system, the evaluation harness, and
the cycle loop all read from it. The record carries the raw answer string kept
verbatim for scoring, the tier that handled the query, a flag for whether the
answer was stored, the full routing decision, the pre-routing confidence, the
complete verifier record with all nine signals and the gate decision, the single
trust scalar, the memory identifier if one was written, an escalation flag, and a
display answer that maps an abstention to a refusal string and a discard to an
empty surface.

\begin{bythenumbers}[The four gate outcomes and their thresholds]
The gate of \Cref{ch:composite} maps the calibrated composite to one of four
outcomes. Holding the locked operating point in mind makes the scenarios above
exact rather than approximate.
\begin{itemize}[leftmargin=1.3em, itemsep=1pt]
  \item \textbf{Store} when the composite is at least $\taustore = 0.60$: write
  the episode into high-precision memory.
  \item \textbf{Defer} when it lands in $[\taudefer, \taustore) = [0.45, 0.60)$:
  park it in the deferred buffer for reconsideration next cycle.
  \item \textbf{Abstain} when it is below $\taudefer$ \emph{and} an abstention
  condition holds (for instance grounding below its floor): return a calibrated
  refusal.
  \item \textbf{Discard} otherwise: drop the answer silently without storing.
\end{itemize}
A single fixed pair of thresholds works across five very different benchmarks
because the signals are \emph{per-benchmark calibrated} before they are compared,
a point we develop in \Cref{ch:composite}.
\end{bythenumbers}

\section{The cycle loop that wraps it all}
\label{sec:cycle-loop}

Now the outer clock. A cycle boundary is not a single step but an ordered block,
and the order matters because some steps depend on the model having actually
improved. \Cref{fig:cycle-loop} sketches the block; the precise mechanics are the
subject of \Cref{ch:cycle}.

\begin{figure}[htbp]
\centering
\begin{tikzpicture}[
  font=\small\sffamily,
  st/.style={rectangle, rounded corners=4pt, line width=1pt, align=center,
    inner sep=5pt, minimum height=1.0cm, text width=3.5cm,
    drop shadow={shadow xshift=0.3ex,shadow yshift=-0.3ex,fill=black!10}},
  commit/.style={st, draw=cbExample!70!black, fill=cbExample!10},
  always/.style={st, draw=cbNumbers!70!black, fill=cbNumbers!10},
  fork/.style={diamond, aspect=2, draw=cbPitfall!75!black, fill=cbPitfall!10,
    line width=1pt, align=center, inner sep=2pt, font=\footnotesize,
    text width=2.2cm},
  flow/.style={-{Stealth[length=2.4mm]}, line width=1.1pt, draw=cbInk!75},
  no/.style={-{Stealth[length=2.4mm]}, line width=1.1pt, draw=cbPitfall!75,
    dashed},
]
  \node[commit] (ft) at (0,2.4)
    {\textbf{1. Fine-tune}\\{\footnotesize LoRA on verified answers}};
  \node[fork] (guard) at (0,0.55)
    {\textbf{Retention}\\\textbf{guard}};
  \node[commit] (recal) at (5.4,0.55)
    {\textbf{2. Recalibrate +}\\\textbf{retro-verify}\\{\footnotesize commit only}};
  \node[always] (grow) at (5.4,-1.7)
    {\textbf{3. Grow memory}\\{\footnotesize fresh stream, always}};
  \node[commit] (keep) at (0,-1.7)
    {\textbf{Keep new adapter}};
  \node[st, draw=cbInk!55, fill=cbInk!6, text width=3.5cm] (revert) at (-5.4,0.55)
    {\textbf{Roll back adapter}\\{\footnotesize memory kept}};

  \draw[flow] (ft) -- (guard);
  \draw[flow] (guard) -- node[right,font=\scriptsize,text=cbExample!55!black]
    {pass} (keep);
  \draw[no] (guard) -- node[above,font=\scriptsize,text=cbPitfall!70!black]
    {fail} (revert);
  \draw[flow] (keep.east) -- ++(1.4,0) |- (recal.west);
  \draw[flow] (recal) -- (grow);
  \draw[flow] (revert.south) |- ($(keep.south)+(0,-0.6)$) -| (grow.south);
\end{tikzpicture}
\caption[The cycle-boundary block]{The cycle boundary. The model is fine-tuned on
its own verified answers, then a retention guard checks it has not regressed on
held-out probes. On a pass (green) the new adapter is kept, the verifier is
recalibrated, and the stored memory is re-scored. On a fail (red) the adapter is
rolled back to the previous cycle, and recalibration and re-verification are
skipped. Either way (orange), the memory grows with the cycle's fresh data. This
asymmetry, weights revert but memory persists, is the core safety property. (Schematic.)}
\label{fig:cycle-loop}
\end{figure}

\subsection*{The commit-versus-abort fork}

The retraining step is the only place the model can get worse, so it is the only
place the system is willing to undo its own work. After fine-tuning, the loop
runs a set of retention probes on held-out material and compares the worst probe
against a pristine baseline anchored once at the very start. If any probe ratio
drops below the retention floor ($\rhomin = 0.93$), the cycle is declared an
\emph{abort}: the freshly fine-tuned adapter is thrown away and the previous
cycle's adapter is restored.

\begin{precise}[Asymmetric rollback, the single most important safety property]
On an abort, the model weights revert, but the episodic memory does \emph{not}.
Everything the pipeline stored this cycle survives; only the parametric update is
undone. This asymmetry is deliberate. Memory growth is monotone and
high-precision by construction (each episode passed the gate), so it is safe to
keep. A weight update, by contrast, is a global change that can quietly erode
unrelated capabilities, which is exactly the catastrophic forgetting of
\Cref{ch:problem}. So the cheap, reversible, local thing (the adapter) is the
thing we are willing to throw away, and the expensive, accumulated, verified
thing (the memory) is the thing we protect. The realised trajectory hit this fork
for real: at cycle four a retention probe fell to $0.886$, below the $0.93$ floor,
so the adapter rolled back to its cycle-three state while the memory carried
forward untouched.
\end{precise}

\subsection*{What carries across a boundary, and what resets}

The asymmetry generalises into a clean rule about state. Four things persist
across cycle boundaries, and several things are rebuilt fresh each cycle.

\begin{description}[leftmargin=0pt, labelsep=0.5em, style=nextline]
  \item[Carries forward.] The \textbf{episodic memory} (always, even on abort);
  the \textbf{adapter weights} (on commit they advance, on abort they revert to
  the last good state); the \textbf{calibration}, namely the temperature and the
  per-signal curves (refreshed on commit, retained on abort); and the
  \textbf{deferred buffer} of borderline answers awaiting reconsideration.
  \item[Resets each cycle.] The \textbf{training pool} is rebuilt fresh from the
  current memory; the cycle consumes a \textbf{fresh, disjoint slice} of
  questions it has never seen; and the optimiser and dataloader are torn down and
  rebuilt. The \textbf{pristine retention baseline} is the one deliberate
  exception: it is anchored once at the start and never re-anchored, because
  re-anchoring it against an already-trained model would hide the very
  regression it exists to catch.
\end{description}

\begin{pitfall}[An abort is not a no-op cycle]
It is easy to assume that an aborted cycle does nothing. It does a great deal.
Only the \emph{commit-gated} half is skipped: the weight update, the
recalibration, the retroactive re-verification, the deferred-buffer
reconsideration, and the memory consolidation. The \emph{unconditional} half runs
regardless: the pipeline still processes the cycle's fresh data and grows the
memory, still runs the held-out evaluation, still checkpoints memory and the
deferred buffer, and still consults the early-stop gate. An abort costs the
system its weight update for that cycle, not its progress.
\end{pitfall}

\section{The map of the book}
\label{sec:book-map}

With the shape in hand, the rest of the book is a guided descent into each box.
\Cref{tab:book-map} is the index to keep a thumb on: every stage and its chapter.

\begin{table}[htbp]
\centering
\caption[The stage-to-chapter map]{Where each stage of \Cref{fig:master-pipeline}
and \Cref{fig:cycle-loop} is opened in full.}
\label{tab:book-map}
\small
\renewcommand{\arraystretch}{1.25}
\begin{tabularx}{\textwidth}{@{}c >{\raggedright\arraybackslash}X l@{}}
\toprule
\textbf{Stage} & \textbf{What it decides} & \textbf{Chapter} \\
\midrule
1  & Pre-routing confidence $\upre$: how likely is a correct self-answer?
   & \Cref{ch:upre} \\
2  & Episodic memory: encode the query and find the nearest stored episode
   & \Cref{ch:memory} \\
3  & Router: fuse similarity and stored quality into a tier choice
   & \Cref{ch:router} \\
3b & Retrieval-utility classifier: will retrieval actually help here?
   & \Cref{ch:ruc} \\
4  & Tier execution: serve from memory, generate, or retrieve and generate
   & \Cref{ch:tiers} \\
5  & Nine-signal verifier: score the generated answer's trustworthiness
   & \Cref{ch:verifier} \\
6  & Calibrated composite and four-outcome gate: store, defer, abstain, discard
   & \Cref{ch:composite} \\
7  & Output contract: return the answer or a calibrated refusal
   & \Cref{ch:output} \\
8  & Cycle boundary: fine-tune, guard, recalibrate, re-verify, grow
   & \Cref{ch:cycle} \\
\bottomrule
\end{tabularx}
\end{table}

Part IV then steps back from the mechanism. \Cref{ch:theory} proves the
data-purity property that \Cref{ch:problem} promised, namely that a verifier only
modestly better than chance still yields a memory pool cleaner than the verifier
itself. \Cref{ch:benchmarks} lays out the five benchmarks, the baselines, and the
hallucination metrics. \Cref{ch:trajectory} walks the realised run from cycle
zero to cycle five, including the cycle-four abort we have already met.
\Cref{ch:ablations} measures what each piece is actually worth by removing it.
The appendices collect every registered constant (\Cref{ch:constants}) and a
glossary of every symbol (\Cref{ch:glossary}).

\begin{defence}[If an examiner asks: ``what is the one-paragraph summary of how it works?'']
\caem{} answers each question through a seven-stage pipeline that first estimates
how confidently it can answer, then routes the query to one of three tiers,
memory lookup, zero-shot generation, or retrieval-augmented generation, picking
the cheapest tier that will do. Generated answers, but not memory hits, pass
through a nine-signal verifier whose calibrated composite gates a four-outcome
decision: store the answer, defer it, refuse, or discard it. That pipeline never
trains. Around it, a cycle loop periodically fine-tunes the model on the answers
it judged trustworthy, under a retention guard that rolls back any update which
degrades held-out performance, while the verified memory it accumulated persists
regardless. Two clocks: one answers, one improves, and a guard keeps improvement
from becoming regression.
\end{defence}

For a single reference that holds the whole machine at once, \Cref{fig:full-pipeline}
reproduces the complete eight-stage pipeline on one page exactly as the main report
draws it, every stage, formula, and threshold in one view. It is dense by design: the
rest of the book is a guided tour of its boxes, and it rewards returning to as the
trajectory and ablation chapters refer back to specific stages.

\input{chapters/fig_full_pipeline}

We now have the map. The next chapter fills in the handful of background concepts,
namely embeddings, calibration, natural-language inference, low-rank adaptation,
and similarity search, that the per-stage chapters lean on, so that the descent
into each box is self-contained.
